\documentclass[11pt]{article}

\usepackage[margin=1in]{geometry}
\usepackage{amsmath,amssymb,amsthm,mathtools,bm}
\usepackage{booktabs}
\usepackage{microtype}
\usepackage{enumitem}
\usepackage[hidelinks]{hyperref}

\newtheorem{theorem}{Theorem}[section]
\newtheorem{proposition}[theorem]{Proposition}
\newtheorem{lemma}[theorem]{Lemma}
\newtheorem{corollary}[theorem]{Corollary}
\theoremstyle{remark}
\newtheorem{remark}[theorem]{Remark}

\newcommand{\R}{\mathbb R}
\newcommand{\Pp}{\mathbb P}
\newcommand{\ip}[2]{\left\langle #1,#2\right\rangle}
\newcommand{\norm}[1]{\left\lVert #1\right\rVert}
\newcommand{\Frob}{\mathrm F}
\newcommand{\one}{\mathbf 1}
\newcommand{\cI}{\mathcal I}
\newcommand{\cE}{\mathcal E}

\title{Sharp Orders of the Minimum Gaussian Softmax Loss\\
and of Its Improvement over the Null Loss}
\author{Technical note}
\date{August 2026}

\begin{document}
\maketitle

\begin{abstract}
Let \(N=\lceil d^\beta\rceil\), where \(\beta>2\) is fixed, put
\(D=d^2\) and \(h=\log N\), and take Zipf frequencies
\(p_i=i^{-\alpha}/H_{N,\alpha}\) with fixed \(\alpha\geq0\).
For the Gaussian softmax associative-memory loss, \(L(0)=h\).
This note determines the sharp order of the optimal improvement
\(\cI^\star=h-L^\star\) for \(0\leq\alpha\leq1\), and the sharp
order of the residual loss \(L^\star\) for \(\alpha>1\).

Below the critical point, the correct loss scale is
\[
 \Gamma_\alpha
 :=\sum_{i=1}^Np_i\min\{h,Dp_i\}
 =h\sum_{p_i\geq h/D}p_i
   +D\sum_{p_i<h/D}p_i^2.
\]
We prove
\[
 h-L^\star=\Theta_{\alpha,\beta}(\Gamma_\alpha),
 \qquad 0\leq\alpha<1.
\]
The threshold \(h/D\), rather than the radius threshold \(D^{-1}\),
is forced by the fact that one query can improve by at most \(h\).
At \(\alpha=1\),
\[
 h-L^\star=\frac{2}{\beta}h+O_\beta(\log h).
\]
For every fixed \(\alpha>1\),
\[
 L^\star
 =\Theta_{\alpha,\beta}\!\left(D^{1-\alpha}h^\alpha\right).
\]
All logarithmic factors are displayed explicitly.  A proof sketch is
given before the detailed arguments.
\end{abstract}

\section{Model and main theorem}
\label{sec:model}

Fix \(\beta>2\) and \(\alpha\geq0\), and set
\[
 N=\lceil d^\beta\rceil,\qquad D=d^2,\qquad h=\log N,
 \qquad
 H_{N,\alpha}=\sum_{k=1}^Nk^{-\alpha},
 \qquad
 p_i=\frac{i^{-\alpha}}{H_{N,\alpha}}.
\]
Let
\[
 u_1,\ldots,u_N,v_1,\ldots,v_N
 \overset{\mathrm{ind}}{\sim}\mathcal N(0,I_d/d).
\]
For \(W\in\R^{d\times d}\), define
\begin{align}
 \ell_i(W)
 &=
 \log\left(1+\sum_{j\neq i}
 \exp\bigl((u_j-u_i)^\top Wv_i\bigr)\right),\\
 L(W)&=\sum_{i=1}^Np_i\ell_i(W),\\
 L^\star&=\inf_{W\in\R^{d\times d}}L(W),
 \qquad
 \cI^\star=L(0)-L^\star=h-L^\star.
 \label{eq:model}
\end{align}
Here \(L(0)=h\), since \(\ell_i(0)=\log N\).

For \(0\leq\alpha<1\), introduce the loss-adapted clipping set
\[
 H_h:=\{i:p_i\geq h/D\},\qquad T_h:=[N]\setminus H_h,
\]
and the deterministic scale
\begin{equation}
 \Gamma_\alpha
 :=h\sum_{i\in H_h}p_i+D\sum_{i\in T_h}p_i^2
 =\sum_{i=1}^Np_i\min\{h,Dp_i\}.
 \label{eq:Gamma-def}
\end{equation}

\begin{theorem}[Sharp loss and improvement orders]
\label{thm:main}
For every fixed \(\alpha\geq0\) and \(\beta>2\), there exist
positive constants \(c_0,c_1\), depending only on
\((\alpha,\beta)\), such that the following statements hold with
probability at least
\[
 1-d^{-c_0}-e^{-c_1N}
\]
for all sufficiently large \(d\).

\begin{enumerate}[label=(\roman*)]
\item If \(0\leq\alpha<1\), then
\begin{equation}
 \boxed{\qquad
 h-L^\star
 =\Theta_{\alpha,\beta}(\Gamma_\alpha),
 \qquad
 L^\star=h-\Theta_{\alpha,\beta}(\Gamma_\alpha).
 \qquad}
 \label{eq:main-subcritical}
\end{equation}
Moreover,
\begin{equation}
 \Gamma_\alpha\asymp_{\alpha,\beta}
 \begin{cases}
 d^{\,2-\beta},
   &0\leq\alpha<\frac12,\\[1mm]
 d^{\,2-\beta}h,
   &\alpha=\frac12,\\[1mm]
 d^{\,2-2\beta(1-\alpha)},
   &\frac12<\alpha<1,\quad\beta(1-\alpha)\geq2,\\[1mm]
 d^{-\frac{(\beta-2)(1-\alpha)}{\alpha}}
 h^{\,2-\frac1\alpha},
   &\frac12<\alpha<1,\quad\beta(1-\alpha)<2.
 \end{cases}
 \label{eq:Gamma-phase}
\end{equation}
Consequently, the requested explicit \(h-\Theta(\cdot)\) statement is
\begin{equation}
 \boxed{
 L^\star=
 \begin{cases}
 h-\Theta_{\alpha,\beta}(d^{\,2-\beta}),
   &0\leq\alpha<\frac12,\\[1mm]
 h-\Theta_{\beta}(d^{\,2-\beta}h),
   &\alpha=\frac12,\\[1mm]
 h-\Theta_{\alpha,\beta}
   (d^{\,2-2\beta(1-\alpha)}),
   &\frac12<\alpha<1,\quad\beta(1-\alpha)\geq2,\\[1mm]
 h-\Theta_{\alpha,\beta}\!\left(
   d^{-\frac{(\beta-2)(1-\alpha)}{\alpha}}
   h^{\,2-\frac1\alpha}\right),
   &\frac12<\alpha<1,\quad\beta(1-\alpha)<2.
 \end{cases}}
 \label{eq:subcritical-explicit}
\end{equation}

\item If \(\alpha=1\), then
\begin{equation}
 \boxed{\qquad
 h-L^\star=\frac2\beta h+O_\beta(\log h),
 \qquad
 L^\star=\left(1-\frac2\beta\right)h+O_\beta(\log h).
 \qquad}
 \label{eq:main-critical}
\end{equation}

\item If \(\alpha>1\), then
\begin{equation}
 \boxed{\qquad
 L^\star
 =\Theta_{\alpha,\beta}\!\left(D^{1-\alpha}h^\alpha\right)
 =\Theta_{\alpha,\beta}\!\left(
 d^{-2(\alpha-1)}(\log N)^\alpha\right).
 \qquad}
 \label{eq:main-supercritical}
\end{equation}
Equivalently,
\[
 h-L^\star
 =h-\Theta_{\alpha,\beta}\!\left(D^{1-\alpha}h^\alpha\right).
\]
\end{enumerate}
\end{theorem}

\begin{remark}[Why the loss clipping level is \(h/D\)]
\label{rem:two-clippings}
The supplied radius note uses
\(\min\{1,Dp_i\}\), hence the threshold \(p_i=D^{-1}\).
That is the correct clipping for \(\norm{W_\star}_{\Frob}\), but not
for the objective value.  In the loss, the linear gain \(Dp_i\) from
query \(i\) saturates when it reaches the null loss \(h\); therefore
the objective-value threshold is \(Dp_i=h\).
In the last line of \eqref{eq:Gamma-phase}, this distinction produces
the genuine factor
\[
 h^{\,2-1/\alpha}.
\]
Thus a bare \(\Theta\) statement in terms of the radius scale would be
false in that regime.
\end{remark}

\begin{remark}[The exact clipping boundary]
If \(1/2<\alpha<1\) and \(\beta(1-\alpha)=2\), then
\[
 Dp_1\longrightarrow1-\alpha<1<h.
\]
Hence \(H_h=\varnothing\) eventually.  Equality belongs to the
non-clipped branch and gives \(\Gamma_\alpha\asymp d^{-2}\), with no
extra logarithm.
\end{remark}

\section{Proof sketch}
\label{sec:sketch}

\paragraph{Step 1: turn minimization into gain maximization.}
Put
\[
 \bar u=\frac1N\sum_{j=1}^Nu_j,\qquad x_j=u_j-\bar u,\qquad
 \Phi(y)=\log\left(\frac1N\sum_{j=1}^Ne^{x_j^\top y}\right).
\]
Then \(\Phi\geq0\), and the centered identity from the supplied note is
\[
 \ell_i(W)=h+\Phi(Wv_i)-x_i^\top Wv_i.
\]
Therefore
\begin{equation}
 \cI(W):=h-L(W)
 =\sum_{i=1}^Np_i
 \bigl(x_i^\top Wv_i-\Phi(Wv_i)\bigr),
 \qquad
 \cI^\star=\sup_W\cI(W).
 \label{eq:gain-representation}
\end{equation}

\paragraph{Step 2: sharp upper bound below one.}
On \(H_h=\{p_i\geq h/D\}\), the deterministic cap gives at most
\[
 h\sum_{i\in H_h}p_i.
\]
On \(T_h\), the centered-gradient bound controls the linear part by
\[
 2\left(\sum_{i\in T_h}p_i^2\right)^{1/2}
 \norm W_{\Frob}.
\]
Retain the free energy only on a still farther deterministic Zipf tail
whose mass tends to one.  The imported diffuse-tail coercivity theorem
then gives, uniformly in \(W\),
\[
 \cI(W)
 \leq
 h\sum_{i\in H_h}p_i+
2\sqrt{S_h}\,\norm W_{\Frob}
-c\psi\!\left(\frac{\norm W_{\Frob}}d\right),
\quad
 S_h=\sum_{i\in T_h}p_i^2,
\]
where \(\psi(t)=\min\{t^2,t\}\).  Optimizing the last two terms gives
\(O(DS_h)\), proving \(\cI^\star\lesssim\Gamma_\alpha\).

\paragraph{Step 3: matching comparison matrix below one.}
Define
\[
 a_i=\min\{h,Dp_i\},\qquad
 S_a=\sum_{i=1}^Na_i u_iv_i^\top.
\]
A fixed-comparator Gaussian-chaos calculation proves
\[
 \ip{\sum_i p_ix_iv_i^\top}{S_a}
 \geq\frac12\Gamma_\alpha,
 \qquad
 \frac1d\sum_i p_i\norm{S_av_i}_2^2
 \leq C\Gamma_\alpha.
\]
All off-target logits \(x_j^\top S_av_i\), \(j\neq i\), are \(o(1)\),
whereas the self logit is \(a_i+o(1)\leq h+o(1)\).
For a sufficiently small fixed \(\eta>0\),
\[
 \sum_i p_i\Phi(\eta S_av_i)
 \leq C\eta^2\Gamma_\alpha+o(\Gamma_\alpha),
\]
while the linear gain is at least
\(\eta\Gamma_\alpha/2\).  Hence
\(\cI^\star\geq\cI(\eta S_a)\gtrsim\Gamma_\alpha\).

\paragraph{Step 4: the endpoint.}
At \(\alpha=1\), the imported head comparator memorizes
\(K_0\asymp D/h^3\) queries and gives
\[
 L^\star
 \leq\left(1-\frac2\beta\right)h+O_\beta(\log h).
\]
For the reverse bound, a deterministic diffuse tail beginning at a
polylogarithmic multiple of \(D/h\) has weighted null loss
\[
 \left(1-\frac2\beta\right)h-O_\beta(\log h).
\]
Uniform free-energy coercivity absorbs its centered linear term, so no
\(W\) can lower the loss below this scale.

\paragraph{Step 5: above one, estimate the residual loss.}
The high-margin capacity is \(K_c\asymp D/h\).  A capacity-scale
head-memory matrix makes the first \(K_c\) losses negligible and pays
only \(O(h)\) on the remaining Zipf mass:
\[
 L^\star\lesssim h\sum_{i>K_c}p_i
 \asymp D^{1-\alpha}h^\alpha.
\]
For the lower bound, the imported one-block coercivity at
\(k\asymp D/h\), combined with the centered-gradient bound, yields
\[
 \sum_{i=k+1}^{2k}\ell_i(W)\gtrsim kh
\]
uniformly in \(W\).  Since \(p_i\gtrsim k^{-\alpha}\) on that block,
\[
 L(W)\gtrsim hk^{1-\alpha}
 \asymp D^{1-\alpha}h^\alpha.
\]

\section{Inputs imported from the supplied note}
\label{sec:imports}

The source is the companion radius note \cite{radius-note}.  Its scope
is \(0\leq\alpha\leq1\).  The following table records every imported
nonstandard input.
\begin{center}
\begin{tabular}{ll}
\toprule
Input used here & Location in the source note\\
\midrule
Existence and uniqueness & Theorem 1.1\\
Centered identity and deterministic cap & equations (2)--(4)\\
Uniform free-energy lower bound & Lemma 4.2\\
Weighted query coercivity & Lemma 4.3 and Corollary 4.4\\
Centered-gradient bound & Lemma 4.6\\
Centered key covariance upper bound & proof of Lemma 4.7, equation (35)\\
Endpoint head comparator & Lemma 6.1 and equation (54)\\
Capacity-block coercivity & Lemma 6.4, equation (60)\\
\bottomrule
\end{tabular}
\end{center}

We restate the exact forms needed below.  With \(\bar u,x_j,\Phi\)
as in Section~\ref{sec:sketch},
\begin{align}
 \ell_i(W)&=h+\Phi(Wv_i)-x_i^\top Wv_i,
 \qquad \Phi(y)\geq0,
 \label{eq:centered}\\
 p_i\bigl(x_i^\top Wv_i-\Phi(Wv_i)\bigr)
 &\leq p_i\min\{h,(x_i^\top Wv_i)_+\}.
 \label{eq:cap}
\end{align}
Let \(\psi(t)=\min\{t^2,t\}\).  With failure probability at most
\(e^{-cN}+e^{-cd}\),
\begin{equation}
 \Phi(y)\geq c\psi\!\left(\frac{\norm y_2}{\sqrt d}\right)
 \qquad\text{for every }y\in\R^d.
 \label{eq:pointwise-free}
\end{equation}
If \(q\) is a fixed deterministic probability vector satisfying
\[
 q_{\max}\leq(Dh^{12})^{-1},
\]
then, with additional failure probability at most \(d^{-30}\),
\begin{equation}
 \sum_iq_i\Phi(Wv_i)
 \geq c\psi\!\left(\frac{\norm W_{\Frob}}d\right)
 \qquad\text{for every }W.
 \label{eq:query-coercivity}
\end{equation}
For every fixed deterministic vector \(b=(b_i)\), with failure at
most \(2e^{-cd}\),
\begin{equation}
 \left\|\sum_i b_ix_iv_i^\top\right\|_{\Frob}
 \leq2\norm b_2.
 \label{eq:centered-gradient}
\end{equation}
Finally, with failure at most \(e^{-cN}\),
\begin{equation}
 \frac1N\sum_{j=1}^Nx_jx_j^\top\preceq\frac{C}{d}I_d.
 \label{eq:key-covariance}
\end{equation}

At \(\alpha=1\), for
\[
 K_0=\lfloor cD/h^3\rfloor,
\]
the source constructs \(W_{\rm head}\) such that
\begin{equation}
 L(W_{\rm head})
 \leq P_{\{K_0+1,\ldots,N\}}(h+C)+N^{-10}
 \label{eq:endpoint-comparator}
\end{equation}
outside an event of probability \(d^{-40}\).
It also proves that, for a sufficiently large fixed \(A=A(\beta)\),
\[
 k=\left\lceil\frac{AD}{h}\right\rceil,\qquad
 B=\{k+1,\ldots,2k\},
 \label{eq:block-def}
\]
satisfy \(2k\leq N\) and, outside an event of probability \(e^{-cd}\),
\begin{equation}
 \frac1k\sum_{i\in B}\ell_i(W)
 \geq c_b\frac{\sqrt h}{d}\norm W_{\Frob}
 \qquad\text{for every }W.
 \label{eq:block-coercivity}
\end{equation}

\section{Zipf evaluation of the subcritical scale}
\label{sec:zipf}

Throughout this section \(0\leq\alpha<1\).  Integral comparison gives
\begin{equation}
 H_{N,\alpha}\asymp_\alpha N^{1-\alpha},
 \qquad
 p_i\asymp_\alpha N^{\alpha-1}i^{-\alpha}.
 \label{eq:zipf-subcritical}
\end{equation}
Let
\[
 m_h:=\max\{i:p_i\geq h/D\},
\]
with \(m_h=0\) if the set is empty.  Then
\begin{equation}
 \Gamma_\alpha
 =h\sum_{i\leq m_h}p_i+D\sum_{i>m_h}p_i^2.
 \label{eq:Gamma-split}
\end{equation}

\begin{lemma}[Evaluation of \(\Gamma_\alpha\)]
\label{lem:Gamma-evaluation}
The four alternatives in \eqref{eq:Gamma-phase} hold.  In every fixed
subcritical regime,
\[
 \Gamma_\alpha=d^{-\rho_{\alpha,\beta}}h^{s_\alpha}
\text{ up to fixed multiplicative constants}
\]
for some \(0<\rho_{\alpha,\beta}<\beta\) and fixed \(s_\alpha\).
In particular, \(\Gamma_\alpha\to0\) polynomially and
\[
 N^{-1+\eta}=o(\Gamma_\alpha)
\]
for every sufficiently small fixed
\(\eta=\eta(\alpha,\beta)>0\).
\end{lemma}

\begin{proof}
If \(0\leq\alpha<1/2\), then
\[
 \sum_{i=1}^Np_i^2\asymp_\alpha N^{-1}.
\]
If \(m_h>0\), then
\[
 \frac{m_h}{N}
 \asymp_{\alpha,\beta}
 \left(\frac{D}{Nh}\right)^{1/\alpha}
 \quad(\alpha>0),
\]
and
\[
 h\sum_{i\leq m_h}p_i
 \asymp h\left(\frac{m_h}{N}\right)^{1-\alpha}.
\]
Dividing this expression by \(D/N\) gives
\[
 \left(\frac{D}{Nh}\right)^{(1-2\alpha)/\alpha}=o(1).
\]
The square-sum after \(m_h=o(N)\) is still of order \(N^{-1}\).
Thus \(\Gamma_\alpha\asymp D/N=d^{2-\beta}\).
For \(\alpha=0\), \(p_i=N^{-1}<h/D\) eventually, and the same
conclusion is immediate.

If \(\alpha=1/2\), then
\[
 \sum_{i=1}^Np_i^2\asymp\frac{h}{N}.
\]
When \(m_h>0\), the head contribution in
\eqref{eq:Gamma-split} is only \(O(D/N)\), while
\[
 D\sum_{i>m_h}p_i^2
 \asymp\frac{D}{N}\log\frac{N}{m_h}
 \asymp\frac{Dh}{N}.
\]
Hence \(\Gamma_{1/2}\asymp d^{2-\beta}h\).

Now suppose \(1/2<\alpha<1\).  Since
\[
 Dp_1=(1-\alpha)d^{\,2-\beta(1-\alpha)}(1+o(1)),
\]
the set \(H_h\) is empty eventually when
\(\beta(1-\alpha)\geq2\).  At equality the displayed expression tends
to \(1-\alpha<h\).  Therefore
\[
 \Gamma_\alpha
 =D\sum_i p_i^2
 \asymp DN^{-2(1-\alpha)}
 =d^{\,2-2\beta(1-\alpha)}.
\]

If \(\beta(1-\alpha)<2\), then \(H_h\neq\varnothing\) eventually and
\begin{equation}
 m_h
 \asymp_{\alpha,\beta}
 \left(\frac{DN^{\alpha-1}}{h}\right)^{1/\alpha}.
 \label{eq:mh}
\end{equation}
Because \(p_{m_h}\asymp h/D\), integral comparison gives
\[
 h\sum_{i\leq m_h}p_i
 \asymp\frac{h^2m_h}{D},
 \qquad
 D\sum_{i>m_h}p_i^2
 \asymp\frac{h^2m_h}{D}.
\]
Substitution of \eqref{eq:mh} yields
\[
 \Gamma_\alpha
 \asymp
 d^{-\frac{(\beta-2)(1-\alpha)}{\alpha}}
 h^{\,2-\frac1\alpha}.
\]
This proves \eqref{eq:Gamma-phase}.  The assertion about
\(\rho_{\alpha,\beta}\) follows by inspecting the four powers of
\(d\); in each case \(0<\rho_{\alpha,\beta}<\beta\).
\end{proof}

For the comparison proof, define
\begin{equation}
 a_i:=\min\{h,Dp_i\},\qquad
 A^2:=\sum_i a_i^2,\qquad
 J:=\sum_i p_ia_i=\Gamma_\alpha,\qquad
 B_a:=\sum_i p_i a_i^2.
 \label{eq:a-def}
\end{equation}
The following elementary consequences will be used repeatedly:
\begin{equation}
 \frac{A^2}{D}\leq J,\qquad
 B_a\leq hJ,\qquad
 \frac{\norm p_2}{\sqrt J}
 \leq d^{-c_{\alpha,\beta}}h^{C_{\alpha,\beta}}.
 \label{eq:aux-arithmetic}
\end{equation}
The first inequality follows termwise from \(a_i\leq Dp_i\), and the
second from \(a_i\leq h\).  The third follows directly by substituting
\eqref{eq:zipf-subcritical} and the four cases of
\eqref{eq:Gamma-phase}; its exponent is strictly negative in every
fixed subcritical regime.

\section{Sharp improvement for \(0\leq\alpha<1\)}
\label{sec:subcritical}

\subsection{Uniform upper bound}

\begin{proposition}[Sharp upper bound on improvement]
\label{prop:sub-upper}
For every fixed \(0\leq\alpha<1\), with probability at least
\(1-d^{-c}-e^{-cN}\),
\[
 \cI(W)\leq C_{\alpha,\beta}\Gamma_\alpha
 \qquad\text{for every }W\in\R^{d\times d}.
\]
\end{proposition}

\begin{proof}
Let \(H_h,T_h\) be as in \eqref{eq:Gamma-def}, and put
\[
 S_h=\sum_{i\in T_h}p_i^2.
\]
Choose the deterministic far tail
\[
 U=\{i:p_i\leq(Dh^{24})^{-1}\},
 \qquad P_U=\sum_{i\in U}p_i.
\]
Equation \eqref{eq:zipf-subcritical} shows that the left endpoint of
\(U\) is \(o(N)\), and therefore \(P_U=1-o(1)\).  In particular,
\(P_U\geq1/2\) for large \(d\).  For
\[
 q_i=\frac{p_i\one_{\{i\in U\}}}{P_U},
\]
we have
\[
 q_{\max}\leq\frac{2}{Dh^{24}}\leq\frac1{Dh^{12}}.
\]
Thus \eqref{eq:query-coercivity} gives
\begin{equation}
 \sum_{i\in U}p_i\Phi(Wv_i)
 \geq c\psi\!\left(\frac{\norm W_{\Frob}}d\right)
 \quad\text{uniformly in }W.
 \label{eq:far-tail-coercivity}
\end{equation}

Use the cap \eqref{eq:cap} on \(H_h\), retain the full linear part on
\(T_h\), and retain the negative free energy only on \(U\subset T_h\).
With \(r=\norm W_{\Frob}\), equations
\eqref{eq:centered-gradient} and \eqref{eq:far-tail-coercivity} give
\begin{align}
 \cI(W)
 &\leq
 h\sum_{i\in H_h}p_i+
 \ip{\sum_{i\in T_h}p_ix_iv_i^\top}{W}
 -\sum_{i\in U}p_i\Phi(Wv_i)\\
 &\leq
 h\sum_{i\in H_h}p_i+
2\sqrt{S_h}\,r
-c\psi(r/d).
 \label{eq:sub-upper-master}
\end{align}
Set \(t=r/d\) and \(A_h=DS_h\).  Then the last two terms are
\[
 2\sqrt{A_h}\,t-c\psi(t).
\]
For \(0\leq t\leq1\), their supremum is at most \(C A_h\).
For \(t\geq1\), they are nonpositive for all sufficiently large \(d\),
because \(A_h\leq\Gamma_\alpha\to0\).  Hence
\[
 \cI(W)
 \leq h\sum_{i\in H_h}p_i+C DS_h
 \leq C\Gamma_\alpha.
\]
Every event used here is uniform before \(W\) is selected.
\end{proof}

\subsection{The log-clipped comparison lemma}

\begin{lemma}[Log-clipped comparison action]
\label{lem:log-clipped}
Fix \(0\leq\alpha<1\), define \(a_i,A,J,B_a\) by
\eqref{eq:a-def}, and set
\[
 S_a=\sum_{i=1}^Na_i u_iv_i^\top.
\]
There is a deterministic \(\varepsilon_d=o(1)\) such that, outside an
event of probability \(d^{-c}+e^{-cN}\),
\begin{align}
 \ip{\sum_i p_ix_iv_i^\top}{S_a}
 &\geq\frac12J,
 \label{eq:comparison-linear}\\
 \frac1d\sum_i p_i\norm{S_av_i}_2^2
 &\leq CJ,
 \label{eq:comparison-action}\\
 \max_i\left|x_i^\top S_av_i-a_i\right|
 +\max_{i\neq j}\left|x_j^\top S_av_i\right|
 &\leq\varepsilon_d.
 \label{eq:comparison-logits}
\end{align}
One may take
\begin{equation}
 \varepsilon_d
 \leq C\left(
 \sqrt{Jh}
 +\frac{h^2}{d}
 +\frac{h^{3/2}}{\sqrt d}
 +\frac{d\sqrt J}{\sqrt N}\right)=o(1).
 \label{eq:epsilon}
\end{equation}
\end{lemma}

\begin{proof}
We give the concentration details because removing the extraneous
polylogarithmic loss from the old radius comparator is essential.

\emph{Linear action.}
First remove centering and write
\[
 G_0=\sum_i p_i u_iv_i^\top.
\]
Expansion gives
\[
 \ip{G_0}{S_a}
 =
 \sum_i p_i a_i\norm{u_i}_2^2\norm{v_i}_2^2
 +\sum_{i\neq k}p_i a_k
 (u_i^\top u_k)(v_i^\top v_k).
\]
The first sum has mean \(J\), and its centered \(L^2\)-norm is at most
\[
 \frac{C}{\sqrt d}
 \left(\sum_i(p_i a_i)^2\right)^{1/2}
 \leq\frac{CJ}{\sqrt d}.
\]
Distinct unordered pairs in the second sum are orthogonal in
\(L^2\), so its \(L^2\)-norm is at most
\[
 C\frac{\norm p_2 A}{d}.
\]
By \eqref{eq:aux-arithmetic}, the ratio of this bound to \(J\) is
polynomially small.  Gaussian-chaos hypercontractivity at a sufficiently
large fixed moment therefore shows
\[
 \ip{G_0}{S_a}\geq\frac34J
\]
outside an event \(d^{-c}\).
Moreover,
\[
 \sum_i p_ix_iv_i^\top
 =G_0-\bar u\left(\sum_i p_iv_i\right)^\top.
\]
The standard bounds
\[
 \norm{\bar u}_2\leq C N^{-1/2},\qquad
 \left\|\sum_i p_iv_i\right\|_2\leq C\norm p_2,\qquad
 \norm{S_a}_{\Frob}\leq CA
\]
and \(N\gg D\), together with \eqref{eq:aux-arithmetic}, make the
centering correction \(o(J)\).  This proves
\eqref{eq:comparison-linear}.

\emph{Quadratic action.}
Let \(U=[u_1,\ldots,u_N]\), \(V=[v_1,\ldots,v_N]\),
\(A_0=\operatorname{diag}(a_i)\), and \(P=\operatorname{diag}(p_i)\).
Conditioning on \(V\),
\[
 Q:=\sum_i p_i\norm{S_av_i}_2^2
 =\operatorname{tr}(UKU^\top),
 \qquad
 K=A_0V^\top VPV^\top VA_0\succeq0.
\]
The same Gaussian quadratic-form estimate used in the source proof of
its equation (34) gives
\begin{equation}
 Q\leq2\operatorname{tr}K
 \label{eq:U-quadratic}
\end{equation}
with conditional failure probability at most \(e^{-cd}\).
Now
\[
 \operatorname{tr}K
 =\sum_{i,k}p_i a_k^2(v_i^\top v_k)^2.
\]
The diagonal part is at most \(CB_a\) on the standard norm event.
For fixed \(k\), condition on \(v_k\).  Weighted Bernstein applied to
\[
 \sum_{i\neq k}p_i(v_i^\top v_k)^2
\]
gives an upper bound \(C/d\).  Indeed, the conditional mean is
\(\norm{v_k}_2^2/d\), the variance proxy is at most
\[
 \frac{C\norm{v_k}_2^4}{d^2}\sum_i p_i^2,
\]
and the largest subexponential scale is at most
\(Cp_1\norm{v_k}_2^2/d\).
Since \(p_1^{-1}\asymp N^{1-\alpha}\) is polynomial in \(d\), the
Bernstein exponent dominates \(\log N\).  A union bound makes
\[
 \max_k\sum_{i\neq k}p_i(v_i^\top v_k)^2\leq\frac Cd
 \label{eq:leave-one-out}
\]
hold with failure \(d^{-c}\).  Consequently,
\[
 \operatorname{tr}K\leq C(B_a+A^2/d).
\]
Combining this with \eqref{eq:U-quadratic} and
\eqref{eq:aux-arithmetic} yields
\[
 \frac Qd
 \leq C\left(\frac{B_a}{d}+\frac{A^2}{D}\right)
 \leq CJ,
\]
which is \eqref{eq:comparison-action}.

\emph{Logits.}
On the standard simultaneous norm and inner-product event, collision
terms contribute at most \(Ca_{\max}\sqrt{h/d}\), where
\(a_{\max}\leq h\).  Conditional Bernstein for the non-collision sum,
followed by a union bound over all ordered pairs, gives
\[
 C\left(\frac{A\sqrt h}{d}+\frac{a_{\max}h}{d}\right).
\]
Finally,
\[
 |\bar u^\top S_av_i|
 \leq C\frac{A}{\sqrt N}.
\]
For the self logit, the diagonal term differs from \(a_i\) by at most
\(Ca_i\sqrt{h/d}\).  Using \(A/d\leq\sqrt J\) gives exactly
\eqref{eq:epsilon}.  Lemma~\ref{lem:Gamma-evaluation} implies that
every term in \eqref{eq:epsilon} tends to zero polynomially, up to
powers of \(h\).  This proves \eqref{eq:comparison-logits}.
\end{proof}

\begin{proposition}[Sharp lower bound on improvement]
\label{prop:sub-lower}
For every fixed \(0\leq\alpha<1\), with probability at least
\(1-d^{-c}-e^{-cN}\),
\[
 \cI^\star\geq c_{\alpha,\beta}\Gamma_\alpha.
\]
\end{proposition}

\begin{proof}
Work on the event of Lemma~\ref{lem:log-clipped} and
\eqref{eq:key-covariance}.  Let
\[
 z_{ji}=x_j^\top S_av_i.
\]
Since \(\sum_jx_j=0\), one has \(\sum_jz_{ji}=0\).
Choose a sufficiently small fixed
\(\eta=\eta(\alpha,\beta)>0\).  By
\eqref{eq:comparison-logits}, for large \(d\),
\[
 |\eta z_{ji}|\leq\frac14\quad(j\neq i),
 \qquad
 z_{ii}\leq h+o(1).
\]
Therefore
\begin{align}
 \frac1N\sum_{j=1}^Ne^{\eta z_{ji}}
 &=
 1+\frac1N\sum_j
 \bigl(e^{\eta z_{ji}}-1-\eta z_{ji}\bigr)\\
 &\leq
 1+N^{-1+\eta+o(1)}
 +\frac{C\eta^2}{N}\sum_{j\neq i}z_{ji}^2.
\end{align}
Using \(\log(1+x)\leq x\) and
\eqref{eq:key-covariance},
\begin{equation}
 \Phi(\eta S_av_i)
 \leq
 N^{-1+\eta+o(1)}
 +C\eta^2\frac{\norm{S_av_i}_2^2}{d}.
 \label{eq:Phi-comparator}
\end{equation}
Lemma~\ref{lem:Gamma-evaluation} permits \(\eta>0\) to be chosen so
small that \(N^{-1+\eta}=o(J)\).  Averaging
\eqref{eq:Phi-comparator} with weights \(p_i\) and using
\eqref{eq:comparison-action} gives
\[
 \sum_i p_i\Phi(\eta S_av_i)
 \leq C\eta^2J+o(J).
\]
On the other hand, \eqref{eq:comparison-linear} gives
\[
 \sum_i p_i x_i^\top(\eta S_a)v_i
 \geq\frac{\eta}{2}J.
\]
Choose \(\eta\) once more so that \(C\eta^2\leq\eta/4\).
Then, for all sufficiently large \(d\),
\[
 \cI(\eta S_a)\geq\frac{\eta}{8}J
 =c_{\alpha,\beta}\Gamma_\alpha.
\]
Taking the supremum over \(W\) proves the proposition.
\end{proof}

Propositions~\ref{prop:sub-upper} and \ref{prop:sub-lower}, together
with Lemma~\ref{lem:Gamma-evaluation}, prove part (i) of
Theorem~\ref{thm:main}.

\section{The critical point \(\alpha=1\)}
\label{sec:critical}

At the endpoint,
\[
 p_i=\frac1{iH_N},
 \qquad
 H_N=h+\gamma+O(N^{-1}).
\]

\subsection{Upper bound on the residual loss}

From \eqref{eq:endpoint-comparator},
\[
 L^\star
 \leq\frac{H_N-H_{K_0}}{H_N}(h+C)+N^{-10}.
\]
Since
\[
 \log K_0
 =\log D-3\log h+O(1)
 =\frac2\beta h-3\log h+O(1),
\]
the harmonic-number expansion gives
\begin{equation}
 L^\star
 \leq\left(1-\frac2\beta\right)h
 +3\log h+O_\beta(1).
 \label{eq:critical-upper}
\end{equation}
Equivalently,
\[
 \cI^\star
 \geq\frac2\beta h-3\log h-O_\beta(1).
\]

\subsection{Lower bound on the residual loss}

Set
\[
 r_d=\left\lceil\frac{Dh^{24}}{H_N}\right\rceil,
 \qquad
 T=\{r_d,\ldots,N\},
 \qquad
 P_T=\sum_{i\in T}p_i,
 \qquad
 S_T=\sum_{i\in T}p_i^2.
\]
Then \(r_d<N\) for large \(d\), and
\[
 \log r_d
 =2\log d+23\log h+o(1)
 =\frac2\beta h+23\log h+o(1).
\]
Consequently,
\begin{equation}
 P_Th
 =\left(1-\frac2\beta\right)h
 -23\log h+O_\beta(1).
 \label{eq:critical-tail-mass}
\end{equation}
In particular, \(P_T\) is bounded below by a positive constant.
Moreover,
\[
 \max_{i\in T}\frac{p_i}{P_T}
 \leq(Dh^{12})^{-1},
 \qquad
 S_T\leq\frac{P_T}{Dh^{24}}
\]
for large \(d\).

Apply \eqref{eq:query-coercivity} to
\(q_i=p_i\one_{\{i\in T\}}/P_T\), and apply
\eqref{eq:centered-gradient} to
\(b_i=p_i\one_{\{i\in T\}}\).  Nonnegativity of losses outside \(T\)
and \eqref{eq:centered} imply, for every \(W\),
\begin{align}
 L(W)
 &\geq
 P_Th+cP_T\psi(t)-\varepsilon t,
 \qquad
 t=\frac{\norm W_{\Frob}}d,\\
 \varepsilon&=2d\sqrt{S_T}
 \leq2\sqrt{P_T}\,h^{-12}.
 \label{eq:critical-master}
\end{align}
For \(t\geq1\), the last two terms are nonnegative for large \(d\).
For \(0\leq t\leq1\),
\[
 cP_Tt^2-\varepsilon t
 \geq-\frac{\varepsilon^2}{4cP_T}
 \geq-Ch^{-24}.
\]
Using \eqref{eq:critical-tail-mass},
\begin{equation}
 L^\star
 \geq
 \left(1-\frac2\beta\right)h
 -23\log h-O_\beta(1).
 \label{eq:critical-lower}
\end{equation}
Equations \eqref{eq:critical-upper} and
\eqref{eq:critical-lower} prove part (ii) of
Theorem~\ref{thm:main}.

\section{The supercritical regime \(\alpha>1\)}
\label{sec:supercritical}

For fixed \(\alpha>1\),
\[
 H_{N,\alpha}\asymp_\alpha1.
\]
If \(K\to\infty\) and \(K=o(N)\), integral comparison gives
\begin{equation}
 \sum_{i>K}p_i\asymp_\alpha K^{1-\alpha},
 \qquad
 Kp_{2K}\asymp_\alpha K^{1-\alpha}.
 \label{eq:super-tail}
\end{equation}

\subsection{Capacity-scale head memory}

\begin{lemma}[Head memory at scale \(D/h\)]
\label{lem:capacity-memory}
There are constants \(c_{\rm mem},c_0,C_0>0\), depending at most on
\(\beta\), such that, if
\[
 K=\left\lfloor c_{\rm mem}\frac{D}{h}\right\rfloor,
 \qquad
 S_K=\sum_{k=1}^Ku_kv_k^\top,
\]
then, with probability at least \(1-d^{-40}\), simultaneously for
every \(i\neq j\),
\begin{align}
 (u_i-u_j)^\top S_Kv_i&\geq c_0,
 &&i\leq K,
 \label{eq:capacity-head}\\
 |(u_j-u_i)^\top S_Kv_i|&\leq C_0,
 &&i>K.
 \label{eq:capacity-tail}
\end{align}
\end{lemma}

\begin{proof}
Let \(\delta=d^{-40}\) and
\(\Lambda=\log(64N^2/\delta)\asymp_\beta h\).
Standard chi-square and conditional Gaussian inner-product tails give
\[
 \max_i\left|
 \norm{u_i}_2^2-1\right|
 \vee
 \max_i\left|\norm{v_i}_2^2-1\right|
 \leq\frac14,
\]
and
\[
 \max_{i\neq j}
 \{|u_i^\top u_j|,|v_i^\top v_j|\}
 \leq C\sqrt{\Lambda/d}
\]
outside an event of probability at most \(\delta/4\).

Fix \(i\neq j\), condition on
\(\sigma(u_i,u_j,v_i)\), and for
\(k\leq K\), \(k\notin\{i,j\}\), set
\[
 Z_k^{(i,j)}
 =((u_i-u_j)^\top u_k)(v_k^\top v_i).
\]
These variables are conditionally independent and centered, with
\[
 \norm{Z_k^{(i,j)}}_{\psi_1}\leq C/d.
\]
Conditional Bernstein and a union bound over the \(N^2\) ordered
pairs give
\[
 \max_{i\neq j}
 \left|
 \sum_{\substack{k\leq K\\k\notin\{i,j\}}}
 Z_k^{(i,j)}
 \right|
 \leq
 C\left(\frac{\sqrt{K\Lambda}}d+\frac{\Lambda}{d}\right)
 \leq C_\beta\sqrt{c_{\rm mem}}+o(1).
\]
Choose \(c_{\rm mem}\) so that this is at most \(1/16\).
For \(i\leq K\), the self term is at least \(1/2\), while the possible
\(k=j\) collision is \(o(1)\), proving
\eqref{eq:capacity-head}.  If \(i>K\), there is no self term, and the
same estimate proves \eqref{eq:capacity-tail}.
\end{proof}

\begin{proposition}[Supercritical upper bound]
\label{prop:super-upper}
For every fixed \(\alpha>1\), with probability at least \(1-d^{-40}\),
\[
 L^\star\leq
 C_{\alpha,\beta}D^{1-\alpha}h^\alpha.
\]
\end{proposition}

\begin{proof}
Choose a fixed \(T_\alpha>0\) such that
\[
 \beta(T_\alpha c_0-1)>2(\alpha-1)+1,
\]
and set \(W_{\rm mem}=T_\alpha hS_K\).
On the event of Lemma~\ref{lem:capacity-memory},
\[
 \ell_i(W_{\rm mem})\leq N^{1-T_\alpha c_0},
 \quad i\leq K,
\]
whereas
\[
 \ell_i(W_{\rm mem})\leq
 \log(1+Ne^{T_\alpha C_0h})
 \leq C_{\alpha,\beta}h,
 \quad i>K.
\]
Thus
\[
 L(W_{\rm mem})
 \leq N^{1-T_\alpha c_0}
 +C_{\alpha,\beta}h\sum_{i>K}p_i.
\]
By \eqref{eq:super-tail} and \(K\asymp D/h\),
\[
 h\sum_{i>K}p_i
 \asymp D^{1-\alpha}h^\alpha.
\]
The choice of \(T_\alpha\) makes the first term negligible relative
to this scale.
\end{proof}

\subsection{Matching block lower bound}

\begin{proposition}[Supercritical lower bound]
\label{prop:super-lower}
For every fixed \(\alpha>1\), with probability at least \(1-e^{-cd}\),
\[
 L(W)\geq
 c_{\alpha,\beta}D^{1-\alpha}h^\alpha
 \qquad\text{for every }W.
\]
\end{proposition}

\begin{proof}
Take \(k,B\) from \eqref{eq:block-def}, and write
\[
 Q_B(W)=\sum_{i\in B}\ell_i(W),
 \qquad r=\norm W_{\Frob}.
\]
The imported block coercivity gives
\begin{equation}
 Q_B(W)\geq ar,
 \qquad
 a=c_b\frac{k\sqrt h}{d}.
 \label{eq:block-first}
\end{equation}
Apply \eqref{eq:centered-gradient} with
\(b_i=\one_{\{i\in B\}}\):
\[
 \left\|\sum_{i\in B}x_iv_i^\top\right\|_{\Frob}
 \leq2\sqrt k.
\]
The centered identity and \(\Phi\geq0\) imply
\begin{equation}
 Q_B(W)\geq kh-2\sqrt k\,r.
 \label{eq:block-second}
\end{equation}
For \(a,b,H>0\),
\[
 \max\{ar,H-br\}\geq H\frac{a}{a+b}.
\]
Using \eqref{eq:block-first}--\eqref{eq:block-second} and \(kh\geq AD\),
\[
 Q_B(W)
 \geq
 kh\,
 \frac{c_b\sqrt{kh}/d}{c_b\sqrt{kh}/d+2}
 \geq c_Akh.
\]
Since \(p_i\geq p_{2k}\) on \(B\),
\[
 L(W)\geq p_{2k}Q_B(W)
 \geq c_{\alpha,\beta}hk^{1-\alpha}
 \asymp D^{1-\alpha}h^\alpha.
\]
\end{proof}

Propositions~\ref{prop:super-upper} and
\ref{prop:super-lower} prove part (iii) of
Theorem~\ref{thm:main}.

\section{Probability ledger and final phase diagram}
\label{sec:probability}

For \(0\leq\alpha<1\), intersect the far-tail coercivity event, the
single deterministic gradient event for \(T_h\), the log-clipped
comparison event, and the centered covariance event.  Their total
failure probability is \(d^{-c}+e^{-cN}\).

At \(\alpha=1\), intersect the endpoint-comparator event, the
far-tail coercivity event, its deterministic gradient event, and the
existence event if one wishes to identify the infimum with the unique
finite minimizer.  The same failure form holds.

For \(\alpha>1\), intersect the capacity-memory event, the imported
block-coercivity event, and the deterministic block-gradient event.
Again the total failure probability is \(d^{-c}+e^{-cN}\).

No estimate treats an adaptive optimizer as independent of the data:
every bound used at an optimizer is deterministic or uniform in \(W\),
and every comparison matrix has deterministic coefficients.

The complete answer is
\[
\boxed{
\begin{array}{c|c|c}
\text{regime}
&h-L^\star
&L^\star\\ \hline
0\leq\alpha<1
&\Theta_{\alpha,\beta}(\Gamma_\alpha)
&h-\Theta_{\alpha,\beta}(\Gamma_\alpha)\\[1mm]
\alpha=1
&\dfrac2\beta h+O_\beta(\log h)
&\left(1-\dfrac2\beta\right)h+O_\beta(\log h)\\[2mm]
\alpha>1
&h-\Theta_{\alpha,\beta}(D^{1-\alpha}h^\alpha)
&\Theta_{\alpha,\beta}(D^{1-\alpha}h^\alpha)
\end{array}}
\]
with \(\Gamma_\alpha\) evaluated explicitly in
\eqref{eq:Gamma-phase}.  In particular, for fixed
\(\alpha<1\), the improvement tends to zero at the displayed sharp
rate; at \(\alpha=1\), it is a fixed fraction of \(h\); and for fixed
\(\alpha>1\), the residual loss tends to zero at the displayed sharp
rate.

\begin{thebibliography}{9}
\bibitem{radius-note}
Technical note,
\emph{Sharp Frobenius Radius of the Gaussian Softmax
Associative-Memory Minimizer for Zipf Frequencies},
July 2026.

\bibitem{tropp}
J.~A. Tropp,
\emph{User-Friendly Tail Bounds for Sums of Random Matrices},
Foundations of Computational Mathematics, 2012.

\bibitem{vershynin}
R.~Vershynin,
\emph{High-Dimensional Probability},
Cambridge University Press, 2018.
\end{thebibliography}

\end{document}
